\documentclass[8pt, aspectratio=169]{beamer}
\usepackage[utf8]{inputenc}
\usepackage[russian]{babel}
\usepackage{amsmath, amssymb, amsfonts}
\usepackage{graphicx}
\usepackage{hyperref}
\usepackage{xcolor}
\usepackage{wrapfig}
\usepackage{float}
\usepackage{gensymb}
\usepackage{textcomp}
\usepackage{multirow}

% Настройка цвета формул
\definecolor{mathcolor}{RGB}{139,69,19} % brown
\definecolor{darkgreen}{RGB}{0,100,0} % темно-зеленый (dark green)
\everymath{\color{mathcolor}}
\everydisplay{\color{mathcolor}}
\linespread{1.2}

% Настройка темы Beamer
\usetheme{Madrid}
\usecolortheme{default}

% Убираем навигационные символы
\setbeamertemplate{navigation symbols}{}

% Определение команд
\newcommand{\Def}{{\color{darkgreen}\boxed{\color{darkgreen}\mathbf{Def:}}}}
\newcommand{\Th}[1]{{\color{darkgreen}\boxed{\color{darkgreen}\mathbf{Th~#1:}}}}
\newcommand{\Ex}[1]{{\color{darkgreen}\boxed{\color{darkgreen}\mathbf{Example~#1:}}}}
\newcommand{\ra}{\rightarrow}
\newcommand{\Ra}{\Rightarrow}
\newcommand{\hra}{\hookrightarrow}

\title{\textbf{2.5 Правила дифференцирования. \newline Дифференцирование сложной и обратной функции. \newline Таблица производных.}}
\author{}
\institute{}
\date{}

\begin{document}

% Титульный слайд
\begin{frame}
\titlepage
\end{frame}

% Слайд 2
\begin{frame}
\frametitle{Производная суммы, произведения и частного}

$\Th{2.5.1}$ \textbf{о связи производной с арифметическими операциями}

Если функции $f$ и $g$ дифференцируемы в точке $x_0$, то

\textbf{1.} линейность дифференцирования:
$(f + g)'(x_0) = f'(x_0) + g'(x_0),$ $\forall C \in \mathbb{R} \hookrightarrow (Cf)'(x_0) = Cf'(x_0);$

\textbf{2.} дифференцирование произведения (правило Лейбница):
$(f \cdot g)'(x_0) = f'(x_0)g(x_0) + f(x_0)g'(x_0);$

\textbf{3.} дифференцирование частного:
если дополнительно $g(x_0) \neq 0$, то

$$\left(\frac{f}{g}\right)' (x_0) = \frac{f'(x_0)g(x_0) - f(x_0)g'(x_0)}{g^2(x_0)}$$
\end{frame}

% Слайд 3
\begin{frame}
\frametitle{Производная суперпозиции}

$\Th{2.5.2}$ \textbf{о производной суперпозиции}

Пусть функция $f$ дифференцируема в т. $x_0$, а функция $g$ дифференцируема в т. $y_0 = f(x_0)$. Тогда в некоторой окрестности т. $x_0$ определена сложная функция
$$h(x) = (g \circ f)(x) := g(f(x)),$$
которая дифференцируема в т. $x_0$ по правилу
$$h'(x_0) = (g \circ f)'(x_0) = g'(f(x_0)) \cdot f'(x_0).$$

Т.е. производная суперпозиции функций равна произведению их производных.
\end{frame}

% Слайд 4
\begin{frame}
\frametitle{Производная обратной функции}

$\Th{2.5.3}$ \textbf{о производной обратной функции}

\textbf{1.} Пусть функция $f$ определена в некоторой окрестности точки $x_0$, строго монотонна и непрерывна в ней.

\textbf{2.} Пусть существует конечная, отличная от нуля производная $f'(x_0) \neq 0$.

Тогда обратная функция дифференцируема в точке $y_0 = f(x_0)$ и $(f^{-1})'(y_0) = 1/f'(x_0)$.

\vspace{0.5cm}

Т.е., производная обратной функции в точке $y_0$ равна числу, которое обратно производной прямой функции, взятой в прообразе $x_0$.
\end{frame}

% Слайд 5-6: Объединённый слайд с таблицами слева направо
\begin{frame}
    \frametitle{Таблица производных}

    \begin{columns}[T, onlytextwidth]

    \begin{column}{0.48\textwidth}
    \centering
    \textbf{Степень, логарифм, экспонента}
    \vspace{0.2cm}

    \begin{tabular}{|c|c|}
    \hline
    Функция $f(x)$ & Производная $f'(x)$ \\
    \hline
    $x^n$ & $n \cdot x^{n-1}$ \\
    \hline
    $x$ & $1$ \\
    \hline
    $|x|$ & $\frac{x}{|x|} = \frac{|x|}{x} = \text{sign}(x)$ при $x \ne 0$ \\
    \hline
    $a^x$ & $a^x \cdot \ln a$ \\
    \hline
    $e^x$ & $e^x$ \\
    \hline
    $\log_a x$ & $\frac{1}{x \cdot \ln a}$ \\
    \hline
    $\ln x$ & $\frac{1}{x}$ \\
    \hline
    \end{tabular}
    \end{column}

    \begin{column}{0.48\textwidth}
    \centering
    \textbf{Тригонометрические и гиперболические}
    \vspace{0.2cm}

    \begin{tabular}{|c|c|}
    \hline
    Функция $f(x)$ & Производная $f'(x)$ \\
    \hline
    $\sin x$ & $\cos x$ \\
    \hline
    $\cos x$ & $-\sin x$ \\
    \hline
    $\operatorname{tg} x$ & $\frac{1}{\cos^2 x}$ \\
    \hline
    $\operatorname{ctg} x$ & $-\frac{1}{\sin^2 x}$ \\
    \hline
    $\arcsin x$ & $\frac{1}{\sqrt{1 - x^2}}$ \\
    \hline
    $\arccos x$ & $-\frac{1}{\sqrt{1 - x^2}}$ \\
    \hline
    $\operatorname{arctg} x$ & $\frac{1}{1 + x^2}$ \\
    \hline
    $\operatorname{arcctg} x$ & $-\frac{1}{1 + x^2}$ \\
    \hline
    $\operatorname{sh} x$ & $\operatorname{ch} x$ \\
    \hline
    $\operatorname{ch} x$ & $\operatorname{sh} x$ \\
    \hline
    $\operatorname{th} x$ & $\frac{1}{\operatorname{ch}^2 x}$ \\
    \hline
    $\operatorname{cth} x$ & $-\frac{1}{\operatorname{sh}^2 x}$ \\
    \hline
    \end{tabular}
    \end{column}

    \end{columns}
\end{frame}

% Слайд 7
\begin{frame}
\frametitle{Примеры}

$\Ex{1}$

Дана функция $y = x^x$, найти производную

$\square$

$$
y' = (x^x)' = (e^{\ln{x}\cdot x})' = (\ln{x} + 1)e^{\ln{x} \cdot x} = (\ln{x} + 1)\cdot x^x
$$

$\blacksquare$

\vspace{0.5cm}

$\Ex{2}$

Найти производную функции $|x|$

$\square$

$$
(|x|)' = (\sqrt{x^2})'=2x \cdot \frac{1}{2}\frac{1}{\sqrt{x^2}} = \frac{x}{|x|}
$$

$$
\frac{x}{|x|} = \frac{|x|}{x} = \text{sign}(x), \; x \ne 0
$$

$\blacksquare$
\end{frame}

% Слайд 8
\begin{frame}
\frametitle{Примеры (продолжение)}

$\Ex{3}$

Найти производную функции $f(x) = \ln{|\sin{x}|}$

$\square$

$$
f'(x) = (\ln{|\sin{x}|})'=\cos{x} \cdot \frac{|\sin{x}|}{\sin{x}} \cdot \frac{1}{|\sin{x}|} = \ctg{x}
$$

$\blacksquare$

\vspace{0.5cm}

$\Ex{4}$

Найти производную функции $f(x) = x^n$ через определение производной

$\square$

$$f'(x) = \lim_{\Delta x \to 0} \frac{f(x + \Delta x) - f(x)}{\Delta x}$$

$$(x^n)' = \lim_{\Delta x \to 0} \frac{(x + \Delta x)^n - x^n}{\Delta x}$$

\end{frame}

% Слайд 9
\begin{frame}
\frametitle{Примеры (продолжение)}

$$(x + \Delta x)^n = x^n + n \cdot x^{n-1} \cdot \Delta x + \frac{n(n-1)}{2} \cdot x^{n-2} \cdot (\Delta x)^2 + \dots + (\Delta x)^n$$

$$(x + \Delta x)^n - x^n = n \cdot x^{n-1} \cdot \Delta x + \frac{n(n-1)}{2} \cdot x^{n-2} \cdot (\Delta x)^2 + \dots + (\Delta x)^n$$

$$\frac{(x + \Delta x)^n - x^n}{\Delta x} = n \cdot x^{n-1} + \frac{n(n-1)}{2} \cdot x^{n-2} \cdot \Delta x + \dots + (\Delta x)^{n-1}$$

$$(x^n)' = \lim_{\Delta x \to 0} \left( n \cdot x^{n-1} + \frac{n(n-1)}{2} \cdot x^{n-2} \cdot \Delta x + \dots + (\Delta x)^{n-1} \right) = n \cdot x^{n-1}$$

$$(x^n)' = n \cdot x^{n-1}$$

$\blacksquare$
\end{frame}

% Слайд 10
\begin{frame}
\frametitle{Примеры (продолжение)}

$\Ex{5}$

Найти производную функции $y = \arcsin x$ через производную обратной функции

$\square$

\textbf{1.} Пусть $y = \arcsin x$, где $x \in [-1; 1]$, а $y \in [-\frac{\pi}{2}; \frac{\pi}{2}]$.
Тогда обратная функция имеет вид:
$$x = \sin y$$

По теореме о производной обратной функции:
$$(\arcsin x)' = \frac{1}{(\sin y)'_y} = \frac{1}{\cos y}$$
где $\cos{y} \ne 0$ при $y = \pm\frac{\pi}{2}n, \; n \in \mathbb{N}$
\end{frame}

% Слайд 11
\begin{frame}
\frametitle{Примеры (продолжение)}

\textbf{2.} Из основного тригонометрического тождества имеем
$$\cos y = \pm \sqrt{1 - \sin^2 y}$$

Но $\cos y \ge 0$ при $y \in [-\frac{\pi}{2}; \frac{\pi}{2}]$,

И помним также, что $\cos{y} \ne 0$ при $y = \pm\frac{\pi}{2}n, \; n \in \mathbb{N}$

Т.е. $\cos y > 0$, следовательно:

$$\cos y = \sqrt{1 - \sin^2 y}$$

\textbf{3.} Поскольку $\sin y = x$, получаем:
$$\cos y = \sqrt{1 - x^2}$$
$$(\arcsin x)' = \frac{1}{\sqrt{1 - x^2}}$$
при  $|x| < 1$

$\blacksquare$
\end{frame}

\end{document}